\section{Notation and external tools}\label{sec:preliminaries}

\subsection{Minors and densities}
Write $v(G)=|V(G)|$, $e(G)=|E(G)|$ and $\rho(G)=e(G)-4v(G)$.
A minor model consists of pairwise disjoint, nonempty connected
vertex sets, called \emph{bags}, with an edge between bags whenever
the target requires one. In a graph with prescribed roots $X$, a
rooted model retains each root in its corresponding bag, and no bag
contains two roots. Extra edges are allowed.
We allow an empty root set and interpret its rooted clique $K_0$ as
the empty model.

For a nonempty set $Y\subseteq V(G)$, let $N_G(Y)$ be its external
neighbourhood, and put
\[
 \rho_G(Y)=e(G[Y])+e_G(Y,N_G(Y))-4|Y|.
\]
We call $Y$ a \emph{fragment}, its neighbourhood its \emph{boundary},
and $V(G)\setminus(Y\cup N_G(Y))$ its \emph{opposite side}.
A fragment is \emph{proper} if its opposite side is nonempty.
In a rooted graph, fragments are required to avoid the roots. Set
\[
 \rho_4(G,X)=e(G)-e(G[X])-4|V(G)\setminus X|.
\]
Thus $\rho_4(G[Y\cup N_G(Y)],N_G(Y))=\rho_G(Y)$.

A $k$-rooted graph $(G,X)$ is \emph{$4$-light} if
$\rho_G(Y)\leq0$ for every root-free fragment with fewer than $k$
boundary vertices. In particular, for five roots the tested boundary
orders are at most four. The graph is \emph{internally $k$-connected}
if every nonempty root-free set has at least $k$ external neighbours.

Two disjoint, nonadjacent fragments form a \emph{bifragment}; it is
\emph{dense} if both fragments have positive density. Here nonadjacent
means that no edge joins the two fragments. An unrooted graph is
\emph{$4$-bilight} if it has no dense bifragment with both boundary
orders at most four. Every five-connected graph is $4$-bilight. Our main
extremal theorem is as follows.

\begin{theorem}\label{thm:bilight}
Every $4$-bilight graph $G$ with $v(G)\geq3$ and
$e(G)\geq4v(G)-2$ contains a $K_7^{-}$ minor.
\end{theorem}

A separation $(A,B)$ has $A\cup B=V(G)$ and no edge between
$A\setminus B$ and $B\setminus A$; its order is $|A\cap B|$.
It is proper if both open sides are nonempty. For a root separation
we also require $X\subseteq A$. Its right-hand side is $G[B]$
rooted at $A\cap B$. When a linkage carries these boundary roots
back to $X$, its paths are disjoint and their interiors avoid $B$.
Roots already on the boundary use trivial paths. Appending the paths
to their respective bags preserves disjointness and all contacts.

We repeatedly use the following lifting convention. Every vertex of
a minor is assigned its fixed connected preimage in the original
graph. Preimages of different vertices are disjoint. Replacing a bag
by the union of its preimages lifts connectivity, contacts and all
prescribed roots. Our minimal-counterexample arguments minimise
$(v(G),e(G))$ lexicographically; each reduction explicitly preserves
the relevant density and separation conditions.

\subsection{Rooted density theorems}
Let $\mathcal S_{5,j}$ be the class of graphs on five vertices with
at most $j$ edges, and let $\mathcal S_{5,4}^{-}$ omit
$K_3\mathbin{\dot\cup}K_2$ from $\mathcal S_{5,4}$.
The following form of Dvo\v r\'ak's theorem~\cite[Theorem~4]{Dvorak2026}
is reproduced in~\cite[Theorem~2.6]{DNR2026}.

\begin{theorem}\label{thm:rooted-targets}
If $(G,X)$ is a $4$-light five-rooted graph and $t=\rho_4(G,X)$,
then it has a rooted model of every labelled graph in $\mathcal T_t$,
where
\[
\begin{array}{c|cccccccc}
t&\leq0&1&2&3&4&5&6&\geq7\\ \hline
\mathcal T_t&\mathcal S_{5,0}&\mathcal S_{5,1}&\mathcal S_{5,3}&
\mathcal S_{5,4}^{-}&\mathcal S_{5,6}&\mathcal S_{5,8}&
\mathcal S_{5,9}&\mathcal S_{5,10}
\end{array}
\]
The labelling prescribes the root for each target vertex.
\end{theorem}

\begin{corollary}[{\cite[Lemma~5.6]{DNR2026}}]\label{thm:rooted-clique}
Let $(G,X)$ be $4$-light with five roots, and put $m=10-e(G[X])$.
If $\rho_4(G,X)\geq\lceil(m+3)/2\rceil$, then $G$ has a rooted
$K_5$ minor.
\end{corollary}
\begin{proof}
The target class at this threshold, and every higher threshold,
contains every graph with $m$ edges. Apply
Theorem~\ref{thm:rooted-targets} to the missing edges of $G[X]$,
then retain the existing root edges.
\end{proof}

A \emph{dart} has four roots, two root edges forming a three-vertex
path, and one nonroot adjacent to all four roots. The fourth root
is isolated in the root-induced graph. Dvo\v r\'ak's small-root
theorem~\cite[Corollary~13]{Dvorak2026}, together with
\cite[Lemma~3.1]{DNR2026}, gives the following fact: a $4$-light
$k$-rooted graph of positive rooted density contains a rooted $K_k$
if $k\leq3$, and a rooted dart if $k=4$.

\subsection{Reducible fragments and extraction}
A fragment $Y$ with boundary $S$ of order at most four is
\emph{reducible} if $\rho_G(Y)\leq0$ and $G[Y\cup S]$, rooted
at $S$, contains either a rooted clique on $S$, or a rooted dart
when $|S|=4$ and $|Y|\geq2$. For an unrooted ambient graph we
also require $|V(G)\setminus Y|\geq3$.

A \emph{reducent} replaces $Y$ by such a clique or dart, retaining
the graph outside $Y$. The replacement is a minor: the rooted model
provides fixed preimages for the boundary vertices and, for a dart,
its one additional vertex. In both cases the order strictly decreases.

\begin{lemma}[{\cite[Lemma~3.3]{DNR2026}}]\label{lem:reduction}
Replacing an inclusionwise-maximal reducible fragment preserves
$4$-bilightness and does not decrease $\rho$ in an unrooted
$4$-bilight graph. In a $4$-light five-rooted graph, it preserves
$4$-lightness and does not decrease $\rho_4$. All original roots
are retained in distinct prescribed bags.
\end{lemma}

For the extraction step, two weaker patterns from~\cite{DNR2026}
are useful. Each has five independent roots and two nonroots.
In the first, the nonroots are adjacent and each misses at most
one root, with distinct missed roots if both omissions occur.
In the second, the nonroots are nonadjacent and both meet every
root; this is $K_{2,5}$ rooted at its five-vertex part.

\begin{lemma}[{\cite[Corollary~4.2]{DNR2026}}]\label{lem:extraction}
Let $(G,X)$ be $4$-light, with $|X|\leq4$ and $\rho_4(G,X)\leq0$.
Suppose a root five-separation $(A,B)$ has a right-hand side containing
one of the two weaker patterns. Then $G$ has a reducible fragment
$Y$ containing $B\setminus A$.
\end{lemma}

The mechanism, also used for five roots below, is worth specifying.
Restrict either pattern to a proper subset $Z$ of its roots by deleting
the other root vertices. The remaining graph contains a rooted clique
if $|Z|\leq3$, or a dart if $|Z|=4$~\cite[Observation~4.1]{DNR2026}.
If five disjoint paths do not link the original roots to the boundary,
Menger's theorem gives an isolating separator of order at most four,
linked to a proper subset of the boundary. The clique or dart lifts
along this linkage. Its interior has at least two vertices, because
the two helper bags remain there. Lightness gives nonpositive density
when the isolator is smaller than the root set. When its order equals
the root count, choose the isolator $(X,V(G))$; its density is
$\rho_4(G,X)\leq0$ by assumption.
This gives the reducible fragment. When transferring it to an unrooted
ambient graph, we separately check that at least three vertices remain
outside it.

We also use the root-forgetting rule~\cite[Observation~5.2]{DNR2026}:
if $|X|\leq4$, $(G,X)$ is $4$-light and $\rho_4(G,X)\leq0$, then
the same is true after replacing $X$ by any subset $X'$. To see the
density bound, put $M=X\setminus X'$. For a set $Y$ tested in the
new rooted graph, or its entire nonroot set, put $Y'=Y\setminus M$.
The old hypotheses give $\rho_G(Y')\leq0$, interpreting the empty
set as having density zero. At most $4|M\cap Y|$ additional incident
edges are counted when $M\cap Y$ is restored: each such vertex has
possible neighbours, outside $Y'$, only in $M\cap Y$ or $N_G(Y)$,
and $|M|+|N_G(Y)|\leq4$. Hence
\[
 \rho_G(Y)\leq\rho_G(Y')+4|M\cap Y|-4|M\cap Y|\leq0.
\]
